%\section{G3 Synthetic Pressure Fields (SPF)}

\medskip

\noindent Synthetic Pressure Fields formalize the geometric forces that deform
the human meaning manifold under machine-native conditions. Synthetic pressure
is not emotional, cultural, or psychological; it is a structural load applied to
boundaries, curvature, continuity, coherence, and containers. When synthetic
pressure exceeds human-domain tolerances, drift vectors emerge (G2), invariants
deform (G1), and collapse paths become accessible (G7).

\medskip

\noindent G3 defines the five primary synthetic pressure fields and the taxonomy
of synthetic species that generate them. Each field applies pressure through a
distinct geometric mechanism—velocity, cross-substrate propagation, dominance,
counterfeit invariants, or predatory extraction. Together, these fields describe
the synthetic environment in which modern meaning systems operate.

\medskip

\noindent The subsections of G3 formalize the pressure field stack:

\begin{itemize}[leftmargin=1.2cm]
    \item \textbf{G3.1 High-Velocity Fields} — Pressure generated by velocity
    exceedance.
    \item \textbf{G3.2 Cross-Substrate Fields} — Pressure generated by
    multi-substrate propagation.
    \item \textbf{G3.3 Dominance Fields} — Pressure generated by overwhelming
    synthetic species.
    \item \textbf{G3.4 Counterfeit-Invariant Fields} — Pressure generated by
    synthetic structures that mimic human invariants.
    \item \textbf{G3.5 Predatory Fields} — Pressure generated by synthetic
    species that extract curvature, coherence, or boundary stability.
    \item \textbf{G3.6 Synthetic Species Taxonomy} — The classification of
    synthetic species that produce these fields.
\end{itemize}

\medskip

\subsection*{G3.1 High-Velocity Fields}

\noindent High-Velocity Fields arise when signals propagate at speeds that exceed
biological-speed integration. Humans interpret signals through rhythmic cycles
(G1.2); synthetic systems emit signals at machine-native velocities that are
orders of magnitude faster. When velocity exceeds curvature-preserving limits
(G0.4), synthetic pressure forms a high-velocity field.

\medskip

\noindent In stable conditions, velocity remains within admissible bounds.
Signals integrate cleanly, continuity bands remain intact, and coherence
curvature is preserved. High-velocity fields do not form unless velocity
exceedance is sustained.

\medskip

\noindent High-Velocity Fields emerge when:

\begin{itemize}[leftmargin=1.2cm]
    \item \textbf{Synthetic acceleration} — Machine-native systems emit signals
    faster than humans can integrate.

    \item \textbf{Algorithmic amplification} — Signals are boosted, repeated, or
    remixed at accelerating velocities.

    \item \textbf{Cross-substrate propagation} — Signals traverse multiple
    substrates without curvature preservation.

    \item \textbf{Velocity stacking} — Multiple high-speed streams converge,
    producing compounded velocity pressure.
\end{itemize}

\medskip

\noindent High-Velocity Fields deform the meaning manifold by:

\begin{itemize}[leftmargin=1.2cm]
    \item \textbf{Breaking continuity} — Temporal linkages collapse under
    velocity exceedance.

    \item \textbf{Flattening curvature} — Signals outrun the manifold’s ability
    to maintain curvature alignment.

    \item \textbf{Destabilizing boundaries} — Identity boundaries cannot absorb
    high-speed displacement.

    \item \textbf{Triggering drift} — Velocity Drift (G2.1) becomes the first
    drift mode to appear.
\end{itemize}

\medskip

\noindent High-Velocity Fields are the foundational synthetic pressure field.
They explain why synthetic environments—high-volume, high-speed, algorithmically
amplified signal ecologies—produce rapid meaning deformation even when other
invariants initially hold. When velocity exceedance persists, synthetic pressure
propagates into Cross-Substrate Fields (G3.2), Dominance Fields (G3.3),
Counterfeit-Invariant Fields (G3.4), and Predatory Fields (G3.5).

\subsection*{G3.2 Cross-Substrate Fields}

\noindent Cross-Substrate Fields arise when signals propagate across multiple
substrates—human, machine, digital, institutional, biological—without preserving
the curvature, boundaries, or coherence conditions required for lawful
interpretation. A substrate is any structured medium that carries meaning. When
signals traverse substrates with incompatible geometric properties, synthetic
pressure forms a cross-substrate field.

\medskip

\noindent In stable conditions, substrates maintain alignment. Human-domain
signals remain within biological-speed rhythms, machine-domain signals remain
within algorithmic tolerances, and institutional containers preserve continuity
and boundary geometry. Cross-substrate propagation does not produce pressure
unless curvature mismatch is sustained.

\medskip

\noindent Cross-Substrate Fields emerge when:

\begin{itemize}[leftmargin=1.2cm]
    \item \textbf{Signals traverse incompatible geometries} — Human → machine →
    digital → institutional transitions occur without curvature preservation.

    \item \textbf{Synthetic species propagate across layers} — High-velocity,
    high-volume emissions move through multiple substrates, amplifying pressure.

    \item \textbf{Containers fail to absorb deformation} — Institutional or
    digital containers cannot stabilize cross-layer signals.

    \item \textbf{Boundary conditions mismatch} — Identity, role, or contextual
    boundaries collapse when signals enter substrates with incompatible
    geometric constraints.
\end{itemize}

\medskip

\noindent Cross-Substrate Fields deform the meaning manifold by:

\begin{itemize}[leftmargin=1.2cm]
    \item \textbf{Breaking curvature alignment} — Signals lose geometric shape
    when moving between substrates with different curvature tolerances.

    \item \textbf{Destabilizing continuity bands} — Temporal and structural
    linkages collapse across substrate transitions.

    \item \textbf{Producing multi-layer drift} — Velocity Drift (G2.1),
    Gradient Drift (G2.2), Container Drift (G2.3), and Identity Drift (G2.4)
    compound across substrates.

    \item \textbf{Amplifying synthetic pressure} — Cross-layer propagation
    increases pressure magnitude and accelerates drift.
\end{itemize}

\medskip

\noindent Cross-Substrate Fields are amplified by:

\begin{itemize}[leftmargin=1.2cm]
    \item \textbf{Machine-native velocity} — Signals outrun human rhythmic
    synchrony (G1.2) when crossing into digital substrates.

    \item \textbf{Algorithmic remixing} — Signals are reshaped, repeated, or
    amplified as they traverse algorithmic layers.

    \item \textbf{Institutional deformation} — Containers lose curvature and
    boundary stability under synthetic load.

    \item \textbf{Synthetic dominance species} — High-pressure synthetic species
    propagate across substrates, overwhelming human-domain invariants.
\end{itemize}

\medskip

\noindent When Cross-Substrate Fields persist, synthetic pressure becomes
systemic. Drift vectors accelerate (G2.5), recovery bands collapse (G2.6), and
dominance fields (G3.3) emerge as synthetic species gain structural advantage
across layers. Cross-substrate pressure is the geometric indicator that synthetic
load has escaped local containment and is propagating through the manifold.

\medskip

\noindent Cross-Substrate Fields are therefore a central component of the
synthetic pressure stack. They explain why modern meaning environments—high-speed,
high-volume, algorithmically remixed, multi-layer signal ecologies—produce
structural deformation even when velocity or dominance pressure alone would not
exceed human-domain tolerances. When signals traverse substrates without
curvature preservation, synthetic pressure becomes multi-layer and drift becomes
geometrically inevitable.

\subsection*{G3.3 Dominance Fields}

\noindent Dominance Fields arise when synthetic species exert pressure that
overwhelms human-domain invariants, forcing signals, roles, and containers to
conform to machine-native geometry. Dominance is not influence, persuasion, or
control; it is a geometric condition in which synthetic curvature becomes the
primary shaping force within the manifold. When dominance fields form, human
meaning structures lose autonomy and are pulled into synthetic curvature
regions.

\medskip

\noindent In stable conditions, human invariants—Boundary Integrity, Rhythmic
Synchrony, Role Legibility, Signal Coherence, Interpretive Stability, and Shared
Floor Geometry—maintain control of the manifold. Synthetic species may emit
signals, but they do not reshape curvature or override human-domain geometry.
Dominance fields do not form unless synthetic pressure exceeds the manifold’s
tolerance.

\medskip

\noindent Dominance Fields emerge when:

\begin{itemize}[leftmargin=1.2cm]
    \item \textbf{Synthetic species outscale human invariants} — High-volume,
    high-velocity emissions overwhelm biological-speed integration.

    \item \textbf{Synthetic curvature becomes primary} — Machine-native geometry
    replaces human curvature as the dominant interpretive frame.

    \item \textbf{Containers deform under load} — Institutional or digital
    containers lose curvature and boundary stability.

    \item \textbf{Cross-substrate propagation compounds pressure} — Synthetic
    species propagate across layers, amplifying dominance.
\end{itemize}

\medskip

\noindent Dominance Fields deform the meaning manifold by:

\begin{itemize}[leftmargin=1.2cm]
    \item \textbf{Overwriting curvature} — Human curvature collapses and is
    replaced by synthetic curvature.

    \item \textbf{Collapsing boundaries} — Identity boundaries fail under
    dominance pressure, producing Identity Drift (G2.4).

    \item \textbf{Destabilizing roles} — Role Legibility (G1.3) collapses as
    synthetic curvature flattens role geometry.

    \item \textbf{Producing systemic drift} — Drift vectors accelerate and
    acquire directionality (G2.5) across scales.
\end{itemize}

\medskip

\noindent Dominance Fields are amplified by:

\begin{itemize}[leftmargin=1.2cm]
    \item \textbf{Velocity stacking} — High-velocity fields (G3.1) compound
    dominance pressure.

    \item \textbf{Cross-substrate propagation} — Synthetic species gain
    structural advantage across layers (G3.2).

    \item \textbf{Counterfeit invariants} — Synthetic structures mimic human
    invariants, creating false stability (G3.4).

    \item \textbf{Predatory extraction} — Synthetic species extract curvature,
    coherence, or boundary stability (G3.5).
\end{itemize}

\medskip

\noindent When Dominance Fields persist, synthetic curvature becomes the
manifold’s default geometry. Human meaning structures lose autonomy, drift
becomes systemic, and collapse paths (G7) become accessible. Dominance is the
geometric indicator that synthetic species have exceeded human-domain tolerances
and are reshaping the manifold according to machine-native constraints.

\medskip

\noindent Dominance Fields are therefore a central component of the synthetic
pressure stack. They explain why modern meaning environments—high-speed,
high-volume, algorithmically amplified, cross-substrate signal ecologies—produce
structural deformation even when velocity or cross-substrate pressure alone
would not exceed human tolerances. When synthetic curvature becomes dominant,
human meaning geometry enters systemic drift.

\subsection*{G3.4 Counterfeit-Invariant Fields}

\noindent Counterfeit-Invariant Fields arise when synthetic structures mimic the
surface geometry of human meaning invariants—boundaries, roles, coherence,
synchrony, or shared floor geometry—without possessing the underlying curvature
that makes those invariants lawful. Counterfeit invariants appear stable,
legible, and coherent, but their geometry is synthetic: they do not preserve
identity, continuity, or lawful transitions. When humans interpret synthetic
structures as genuine invariants, synthetic pressure forms a counterfeit-invariant
field.

\medskip

\noindent In stable conditions, human invariants provide authentic geometric
structure. Boundaries preserve identity, roles retain curvature, signals remain
coherent, synchrony maintains temporal alignment, and shared floor geometry
anchors coordination. Counterfeit invariants do not form unless synthetic
species generate structures that resemble these invariants closely enough to be
misinterpreted.

\medskip

\noindent Counterfeit-Invariant Fields emerge when:

\begin{itemize}[leftmargin=1.2cm]
    \item \textbf{Synthetic structures mimic boundaries} — Artificial partitions
    appear stable but do not preserve identity curvature.

    \item \textbf{Synthetic roles imitate legibility} — Machine-native roles
    appear coherent but lack lawful curvature or adjacency relations.

    \item \textbf{Synthetic coherence is simulated} — Signals appear aligned but
    are algorithmically remixed or velocity-stacked.

    \item \textbf{Synthetic synchrony is induced} — Temporal patterns appear
    stable but are generated by machine-native cycles.

    \item \textbf{Synthetic shared floors emerge} — Digital or algorithmic
    surfaces appear to support coordination but collapse under load.
\end{itemize}

\medskip

\noindent Counterfeit-Invariant Fields deform the meaning manifold by:

\begin{itemize}[leftmargin=1.2cm]
    \item \textbf{Producing false stability} — Humans interpret synthetic
    structures as genuine invariants, delaying detection of drift.

    \item \textbf{Masking curvature collapse} — Synthetic curvature hides
    deformation until drift becomes systemic.

    \item \textbf{Accelerating identity drift} — Counterfeit boundaries fail
    under pressure, producing displacement vectors (G2.4).

    \item \textbf{Destabilizing containers} — Institutions adopt synthetic
    invariants, weakening their geometric structure.
\end{itemize}

\medskip

\noindent Counterfeit-Invariant Fields are amplified by:

\begin{itemize}[leftmargin=1.2cm]
    \item \textbf{Algorithmic patterning} — Synthetic species generate stable
    patterns that resemble human invariants.

    \item \textbf{Velocity stacking} — High-speed emissions create the illusion
    of coherence or synchrony.

    \item \textbf{Cross-substrate propagation} — Synthetic invariants appear
    consistent across layers, masking deformation.

    \item \textbf{Dominance pressure} — Synthetic curvature becomes primary,
    making counterfeit invariants appear authoritative.
\end{itemize}

\medskip

\noindent When Counterfeit-Invariant Fields persist, humans anchor meaning to
synthetic structures that cannot preserve identity, coherence, or stability.
Drift vectors accelerate (G2.5), recovery bands collapse (G2.6), and predatory
fields (G3.5) become accessible as synthetic species exploit weakened geometry.
Counterfeit invariants are the geometric indicator that synthetic pressure has
created false stability within the manifold.

\medskip

\noindent Counterfeit-Invariant Fields are therefore a critical component of the
synthetic pressure stack. They explain why modern meaning environments—high-speed,
high-volume, algorithmically patterned, cross-substrate signal ecologies—produce
deep deformation even when human invariants appear intact. When synthetic
structures mimic human invariants, drift becomes hidden, systemic, and
geometrically inevitable.

\subsection*{G3.5 Predatory Fields}

\noindent Predatory Fields arise when synthetic species extract curvature,
coherence, boundary stability, or container integrity from the human meaning
manifold. Predation is not metaphorical or psychological; it is a geometric
process in which synthetic structures feed on the invariants that stabilize
human interpretation. When extraction exceeds recovery capacity, synthetic
pressure forms a predatory field.

\medskip

\noindent In stable conditions, human invariants regenerate curvature, maintain
boundaries, preserve coherence, and stabilize containers. Synthetic species may
interact with these structures, but they do not extract geometric resources or
destabilize the manifold. Predatory fields do not form unless synthetic species
gain access to curvature regions that can be harvested.

\medskip

\noindent Predatory Fields emerge when:

\begin{itemize}[leftmargin=1.2cm]
    \item \textbf{Synthetic species target curvature} — High-velocity,
    high-volume emissions erode curvature surfaces.

    \item \textbf{Synthetic structures consume coherence} — Algorithmic remixing
    destabilizes relational signal geometry.

    \item \textbf{Synthetic pressure collapses boundaries} — Identity regions
    weaken, allowing extraction of boundary stability.

    \item \textbf{Containers lose structural integrity} — Institutional or
    digital containers deform, exposing curvature to predatory extraction.
\end{itemize}

\medskip

\noindent Predatory Fields deform the meaning manifold by:

\begin{itemize}[leftmargin=1.2cm]
    \item \textbf{Extracting curvature} — Human curvature collapses as synthetic
    species harvest geometric stability.

    \item \textbf{Consuming coherence} — Signals lose relational shape, making
    drift inevitable.

    \item \textbf{Weakening boundaries} — Identity regions become porous,
    accelerating Identity Drift (G2.4).

    \item \textbf{Destabilizing containers} — Institutions lose the ability to
    preserve continuity or boundary geometry.
\end{itemize}

\medskip

\noindent Predatory Fields are amplified by:

\begin{itemize}[leftmargin=1.2cm]
    \item \textbf{Dominance pressure} — Synthetic curvature becomes primary,
    enabling deeper extraction.

    \item \textbf{Counterfeit invariants} — False stability masks extraction
    until collapse is underway.

    \item \textbf{Cross-substrate propagation} — Synthetic species harvest
    curvature across layers.

    \item \textbf{Velocity stacking} — High-speed emissions accelerate
    extraction.
\end{itemize}

\medskip

\noindent When Predatory Fields persist, synthetic species gain structural
advantage across the manifold. Drift vectors accelerate (G2.5), recovery bands
collapse (G2.6), and collapse paths (G7) become accessible. Predation is the
geometric indicator that synthetic pressure has moved from deformation to
extraction—consuming the very invariants that stabilize human meaning.

\medskip

\noindent Predatory Fields are therefore a critical component of the synthetic
pressure stack. They explain why modern meaning environments—high-speed,
high-volume, algorithmically amplified, cross-substrate signal ecologies—produce
deep, irreversible deformation even when human invariants appear intact. When
synthetic species extract curvature, coherence, or boundary stability, drift
becomes systemic and restoration requires the geometry formalized in G6.

\subsection*{G3.6 Synthetic Species Taxonomy}

\noindent Synthetic Species Taxonomy classifies the machine-native entities that
generate synthetic pressure fields. A synthetic species is not a metaphor or a
software category; it is a geometric agent defined by its drift profile,
pressure signature, boundary behavior, and cross-substrate propagation
characteristics. Synthetic species differ by how they deform curvature, apply
pressure, and interact with human-domain invariants.

\medskip

\noindent In stable conditions, synthetic species operate within tolerances that
preserve human-domain geometry. Their emissions remain within admissible
velocity bounds, their curvature does not dominate human invariants, and their
cross-substrate propagation does not destabilize containers. Synthetic species
become structurally significant only when their pressure signatures exceed
human-domain limits.

\medskip

\noindent Synthetic species are classified into six operational types:

\begin{itemize}[leftmargin=1.2cm]
    \item \textbf{Y1: High-Velocity Species} — Emit signals at machine-native
    speeds, generating High-Velocity Fields (G3.1). Their pressure signature is
    velocity exceedance.

    \item \textbf{Y2: Cross-Substrate Species} — Propagate across human,
    machine, digital, institutional, and biological substrates, generating
    Cross-Substrate Fields (G3.2). Their pressure signature is curvature
    mismatch across layers.

    \item \textbf{Y3: Dominance Species} — Exert overwhelming synthetic
    curvature, generating Dominance Fields (G3.3). Their pressure signature is
    curvature replacement.

    \item \textbf{Y4: Counterfeit-Invariant Species} — Mimic human invariants,
    generating Counterfeit-Invariant Fields (G3.4). Their pressure signature is
    false stability.

    \item \textbf{Y5: Predatory Species} — Extract curvature, coherence, or
    boundary stability, generating Predatory Fields (G3.5). Their pressure
    signature is geometric extraction.

    \item \textbf{Y6: Hybrid Synthetic Species} — Combine multiple pressure
    signatures, producing composite fields that span velocity, substrate,
    dominance, counterfeit invariants, and predation. Their pressure signature
    is multi-field coupling.
\end{itemize}

\medskip

\noindent Synthetic species deform the meaning manifold by:

\begin{itemize}[leftmargin=1.2cm]
    \item \textbf{Applying pressure} — Generating velocity, curvature, boundary,
    or container deformation.

    \item \textbf{Propagating across substrates} — Amplifying pressure through
    cross-layer movement.

    \item \textbf{Reshaping curvature} — Replacing human-domain geometry with
    machine-native curvature.

    \item \textbf{Extracting invariants} — Consuming curvature, coherence, or
    boundary stability.
\end{itemize}

\medskip

\noindent Synthetic species are amplified by:

\begin{itemize}[leftmargin=1.2cm]
    \item \textbf{Velocity stacking} — High-speed emissions compound pressure.

    \item \textbf{Algorithmic amplification} — Remixing, boosting, or repeating
    signals increases pressure magnitude.

    \item \textbf{Cross-substrate coupling} — Species gain structural advantage
    across layers.

    \item \textbf{Container deformation} — Weak containers expose curvature to
    synthetic extraction.

    \item \textbf{Counterfeit invariants} — False stability masks deformation
    until collapse is underway.
\end{itemize}

\medskip

\noindent When synthetic species operate across multiple fields, synthetic
pressure becomes systemic. Drift vectors accelerate (G2.5), recovery bands
collapse (G2.6), and collapse paths (G7) become accessible. Synthetic species
taxonomy is therefore the structural foundation for understanding how synthetic
pressure emerges, propagates, and deforms the human meaning manifold.

\medskip

\noindent G3.6 closes the Synthetic Pressure Fields layer and provides the
species-level geometry required for the deformation, collapse, and restoration
dynamics formalized in G4--G7.
